\chapter{CONSIDERAÇÕES FINAIS}

Este trabalho investigou o comportamento dos algoritmos TCP CUBIC e TCP BBR e desenvolveu uma variante do CUBIC denominada TCP CUBIC-RTT. A proposta foi implementada no simulador \texttt{ns-3} e preservou a estrutura geral do algoritmo original, incorporando informações do menor RTT observado para adaptar o fator de crescimento cúbico $C$ e o fator de redução multiplicativa $\beta$. Os coeficientes $\alpha$ e $\gamma$ controlam a intensidade dessas alterações.

A implementação foi organizada entre o modelo do protocolo e o programa experimental. Foram instrumentados rastros de janela de congestionamento, limiar de congestionamento, RTT, RTO e dados em trânsito, além da coleta das métricas de fluxo pelo FlowMonitor. Durante o desenvolvimento, foi necessário corrigir a topologia para utilizar um gargalo realmente compartilhado, ajustar o caminho de conexão de um callback de rastreamento e filtrar os fluxos reversos de ACK no processamento dos resultados. Essas decisões foram importantes para que a avaliação representasse a disputa pelo enlace e para evitar distorções na consolidação das métricas.

A avaliação reuniu três campanhas, totalizando 735 simulações: a comparação inicial entre CUBIC, BBR e CUBIC-RTT em três atrasos de gargalo, a campanha de robustez, que ampliou os cenários e variou os coeficientes $\alpha$ e $\gamma$, e a campanha de coexistência, que colocou quatro fluxos de cada variante em competição direta pelo mesmo gargalo.

Na comparação inicial, CUBIC e CUBIC-RTT mantiveram vazão agregada praticamente equivalente, e a proposta teve jitter menor nos três atrasos, ao custo de atraso médio maior e, nos cenários de 50 ms e 100 ms, perda também maior. A campanha de robustez confirmou que o efeito da adaptação depende da condição de rede e da combinação dos coeficientes: nenhuma combinação de $\alpha$ e $\gamma$ dominou o CUBIC original em todos os critérios.

A campanha de coexistência revelou um efeito que não apareceria em execuções separadas de cada algoritmo. Com o aumento do atraso do gargalo, a vazão média por fluxo do CUBIC-RTT passou de 1,2345 Mbps em 10 ms para 1,3746 Mbps em 100 ms, enquanto a do CUBIC caiu de 1,2492 Mbps para 1,1070 Mbps, com a vazão total da rede estável em torno de 9,93 Mbps. A vantagem da variante adaptativa foi, portanto, uma redistribuição da capacidade em seu favor, e não um aumento da capacidade utilizada.

Voltando à pergunta que motivou o trabalho, a resposta tem duas partes. Primeiro, foi possível incorporar informações de RTT à estrutura do CUBIC sem trocar sua função cúbica por outro modelo de controle, e essa incorporação mudou o comportamento observado nas simulações. Segundo, não se confirmou que a mudança melhore ao mesmo tempo vazão, atraso, perda e justiça em todas as condições. A campanha de coexistência foi a mais reveladora nesse ponto: uma variante pode ganhar vantagem competitiva sobre fluxos do algoritmo original mesmo quando a utilização total do gargalo praticamente não muda.

A principal contribuição do trabalho é, portanto, a implementação e a avaliação experimental de uma adaptação do CUBIC baseada em RTT, acompanhada da identificação de seus benefícios e de seus limites. Os resultados não sustentam afirmar que o CUBIC-RTT seja universalmente superior ao CUBIC ou ao BBR. Sustentam, porém, que a informação de RTT é um sinal relevante para modificar a dinâmica do algoritmo e que sua utilização precisa ser avaliada também sob competição com fluxos não modificados.

Como trabalhos futuros, recomenda-se aprofundar a análise do cenário \texttt{access100}, ampliar a avaliação de justiça na coexistência e testar a proposta em outras combinações de capacidade, atraso, quantidade de fluxos e disciplinas de fila. Também é relevante investigar mecanismos que utilizem variações recentes do RTT, em vez de depender apenas do menor valor observado. Durante o desenvolvimento foi considerada uma alternativa baseada na variação relativa do RTT suavizada por uma média móvel exponencial (EWMA), aplicada ao fator $C$. Essa alternativa não integra os resultados apresentados e deve ser tratada como uma hipótese para experimentos posteriores. Por fim, a execução em ambiente real e a avaliação de interações com outros algoritmos de controle de congestionamento poderão ampliar a validade externa dos resultados obtidos no simulador.

\appendix
